\documentclass{amsart}
\usepackage{amsmath,amssymb,amsthm,hyperref}
\usepackage{algorithm}
\usepackage{algpseudocode}

\newtheorem{theorem}{Theorem}
\newtheorem{lemma}[theorem]{Lemma}
\newtheorem{proposition}[theorem]{Proposition}
\newtheorem{corollary}[theorem]{Corollary}
\theoremstyle{definition}
\newtheorem{definition}[theorem]{Definition}
\newtheorem{remark}[theorem]{Remark}

\begin{document}

\title{The Drop-The-Larger-Hopeless policy is not constant-competitive:
$R=\infty$}
\maketitle

\begin{abstract}
We determine the worst-case ratio
$R=\sup_\omega \mathrm{Miss}_{\mathrm{DTLH}}(\omega)/\mathrm{Miss}_{\mathrm{OPT}}(\omega)$
for the Drop-The-Larger-Hopeless (DTLH) online deadline-scheduling policy
against the clairvoyant offline optimum. After formalizing the model we prove
$R=\infty$: DTLH is \emph{not} constant-competitive. The witness is an explicit,
elementary one-parameter family $F_r$ of $r+1$ jobs sharing a single
release/deadline window, on which $\mathrm{Miss}_{\mathrm{OPT}}=1$ while a legal
DTLH execution drops $r-1$ jobs; letting $r\to\infty$ forces the ratio to
infinity. We also record two positive results (DTLH never drops a job when the
input is schedulable; the evict-largest deterministic variant is exactly optimal
on a single window) and discuss the deterministic case.
\end{abstract}

\section{Formalization of the model}

We make the informal statement precise.

\begin{definition}[Instance]
An \emph{instance} is a finite multiset
$\omega=\{(a_j,s_j,d_j)\}_{j=1}^{n}$ of jobs, where $a_j\in\mathbb{R}_{\ge 0}$
is the arrival (release) time, $s_j>0$ the size, and $d_j\ge a_j+s_j$ the
deadline of job $j$. A single server processes work at unit rate and may
preempt at no cost.
\end{definition}

\begin{definition}[Feasibility]\label{def:feas}
Fix a time $t$ and a finite set $Q$ of jobs, each job $j\in Q$ carrying a
\emph{remaining size} $r_j>0$ and a deadline $d_j\ge t$. We say $Q$ is
\emph{feasible at time $t$} if there is a unit-rate single-server schedule of
$Q$ on $[t,\infty)$ completing every $j\in Q$ by $d_j$. Equivalently, by the
exchange (earliest-deadline-first) argument, $Q$ is feasible at $t$ iff, listing
$Q=\{j_1,\dots,j_m\}$ with $d_{j_1}\le\cdots\le d_{j_m}$,
\begin{equation}\label{eq:edf}
t+\sum_{i=1}^{\ell} r_{j_i}\ \le\ d_{j_\ell}\qquad\text{for every }\ell=1,\dots,m.
\end{equation}
A job is \emph{hopeless} at $t$ if it cannot be completed by its deadline under
any continuation, i.e.\ already $t+r_j>d_j$.
\end{definition}

Throughout, the server is allowed to realize \emph{any} feasible schedule of its
currently admitted set; ``feasible'' always refers to
Definition~\ref{def:feas}. (The phrase ``serves in FCFS order'' in the prompt is
only the default tie-break used by DTLH to decide \emph{which jobs are present};
the admission test and the completion guarantee use the feasibility relation
\eqref{eq:edf}, as the prompt itself states: a set is feasible if there
\emph{exists} a schedule meeting all deadlines. Since the admitted set is kept
feasible at every admission epoch, an EDF server completes all admitted jobs on
time, by Horn's/Jackson's theorem on the optimality of EDF for single-machine
preemptive feasibility with release times; the FCFS phrasing is immaterial to
the miss count.)

\begin{definition}[The DTLH policy]\label{def:dtlh}
DTLH maintains an \emph{admitted set} $A$ that is feasible at all times; an
admitted job is guaranteed completion, a non-admitted job is \emph{dropped} and
counts as one miss. Decisions are made online at arrival epochs. When job $j$
arrives at time $t=a_j$, let $A^-$ be the set of admitted, not-yet-completed
jobs at $t$ (with their remaining sizes). DTLH acts as in
Algorithm~\ref{alg:dtlh}.
\end{definition}

\begin{remark}[The size used in the eviction test]\label{rem:sizeconv}
The prompt says DTLH ``compares the new job's size with the \emph{sizes} of
currently queued jobs,'' where size $=$ service requirement $s_j$. We take the
eviction comparison in line~4 to use the jobs' \emph{original} sizes $s_k,s_j$.
For our main theorem (Section~\ref{sec:main}) this distinction is irrelevant:
the witness family lives in a single common window, where no job is ever served
before all admission decisions are made, so original size $=$ remaining size at
every decision. We adopt the original-size reading throughout because it is the
literal meaning of ``size $=$ service requirement'' and because ``size'' is a
fixed attribute known at admission. The \emph{feasibility} test (lines~1 and~4)
necessarily uses remaining work $r_k$, since it concerns whether the surviving
set can still be completed.
\end{remark}

\begin{algorithm}[h]
\caption{DTLH admission decision at the arrival of job $j$ at time $t$}
\label{alg:dtlh}
\begin{algorithmic}[1]
\Require current admitted, not-yet-completed set $A^-$ (remaining sizes $r_k$, original sizes $s_k$), arriving job $j=(s_j,d_j)$, time $t$
\Ensure updated admitted set and the miss count increment
\If{$A^-\cup\{j\}$ is feasible at $t$ (test \eqref{eq:edf} on remaining sizes)}
  \State admit $j$: \ $A^-\gets A^-\cup\{j\}$
\Else
  \State $L\gets\{k\in A^- : s_k> s_j \text{ and } (A^-\setminus\{k\})\cup\{j\}\text{ is feasible at }t\}$
  \If{$L\neq\varnothing$}
    \State pick any $k^\star\in L$; drop $k^\star$, admit $j$: \ $A^-\gets(A^-\setminus\{k^\star\})\cup\{j\}$; one miss ($k^\star$)
  \Else
    \State drop $j$; one miss ($j$)
  \EndIf
\EndIf
\State \Return $A^-$
\end{algorithmic}
\end{algorithm}

\begin{remark}[The policy is underspecified; we measure its worst case]
\label{rem:tiebreak}
Algorithm~\ref{alg:dtlh} contains two genuine sources of nondeterminism: the
order in which simultaneous arrivals (same $a_j$) are processed, and the choice
of victim $k^\star\in L$ in line~5. The prompt fixes neither --- it says
``drops \emph{such} a larger queued job.'' We therefore regard DTLH as a
\emph{family} of policies indexed by these tie-breaks, and define
\[
\mathrm{Miss}_{\mathrm{DTLH}}(\omega)
=\max\{\text{number of jobs dropped over all legal DTLH tie-break choices on }\omega\},
\]
the \emph{worst-case} behaviour of the rule. This is the correct quantity for
the competitive ratio of a policy that is specified only up to tie-breaks: a
single legal execution achieving many misses already witnesses that the rule, as
stated, is not better than the ratio so obtained. Every lower-bound instance
below is presented together with the explicit legal execution realizing the
stated miss count, so all bounds hold for this worst-case reading (and a
fortiori for the supremum over the policy class). The deterministic case is
discussed in Section~\ref{sec:det}.
\end{remark}

\begin{definition}[Benchmarks]
For an instance $\omega$, $\mathrm{Miss}_{\mathrm{DTLH}}(\omega)$ is as in
Remark~\ref{rem:tiebreak}. Let
\begin{equation}\label{eq:opt}
\mathrm{Miss}_{\mathrm{OPT}}(\omega)
=\min\Big\{\,|\omega\setminus S|\ :\ S\subseteq\omega \text{ is feasible respecting release times}\,\Big\},
\end{equation}
the minimum number of jobs that any policy must miss; $S$ is feasible respecting
release times iff its preemptive single-machine EDF schedule completes every job
by its deadline (EDF optimality for fixed job sets with release times,
Horn/Jackson). This $\mathrm{Miss}_{\mathrm{OPT}}$ is the clairvoyant offline
optimum, a lower bound on the misses of every (online or offline) admissible
policy; using it in the denominator \emph{upper-bounds} the achievable
competitive ratio, the strongest way to test optimality of DTLH.
\end{definition}

The headline quantity is
\[
R \;=\; \sup_{\omega:\ \mathrm{Miss}_{\mathrm{OPT}}(\omega)>0}
\frac{\mathrm{Miss}_{\mathrm{DTLH}}(\omega)}{\mathrm{Miss}_{\mathrm{OPT}}(\omega)} .
\]

\section{Two positive structural facts}

\begin{lemma}[No needless drops]\label{lem:nowaste}
If $\omega$ is feasible respecting release times, i.e.\
$\mathrm{Miss}_{\mathrm{OPT}}(\omega)=0$, then \emph{every} legal DTLH execution
drops nothing: $\mathrm{Miss}_{\mathrm{DTLH}}(\omega)=0$.
\end{lemma}

\begin{proof}
We show by induction over arrival epochs that DTLH's admitted set equals the set
of all already-arrived, not-yet-completed jobs (no drop has occurred) and is
feasible at the current time. Initially $A=\varnothing$. Suppose at the arrival
of job $j$ at $t=a_j$ the admitted, not-yet-completed set $A^-$ consists of
exactly the arrived-but-uncompleted jobs. Because $\omega$ is feasible
respecting release times, the EDF schedule of $\omega$ completes all of $\omega$;
restricting that schedule to $[t,\infty)$ yields a feasible completion at time
$t$ of all jobs that have arrived by $t$ and are not yet finished \emph{together
with} $j$ (each such job is in $\omega$, has deadline $\ge t$, and meets its
deadline in the global EDF schedule). Hence $A^-\cup\{j\}$ is feasible at $t$, so
DTLH executes line~2 and admits $j$ without any drop, independently of any
tie-break. The invariant is preserved, and by induction no drop ever occurs.
\end{proof}

Lemma~\ref{lem:nowaste} shows the ratio is nontrivial only when
$\mathrm{Miss}_{\mathrm{OPT}}(\omega)\ge 1$; it also rules out the degenerate
``divide by zero'' regime.

\begin{lemma}[Single common window: the \emph{evict-largest} rule is optimal]
\label{lem:onewindow}
Suppose all jobs share one release time $a$ and one deadline $d$, so feasibility
of a set $T$ is the single constraint $a+\sum_{j\in T}s_j\le d$, i.e.\ total size
at most the capacity $C:=d-a$. Consider the deterministic DTLH variant
$\mathrm{DTLH}_{\max}$ that, on every swap, evicts a victim of \emph{largest}
original size among the eligible set $L$, processing simultaneous arrivals in any
fixed order. Then $\mathrm{DTLH}_{\max}$ keeps a maximum-cardinality feasible
subset, so its miss count equals $\mathrm{Miss}_{\mathrm{OPT}}(\omega)$.
\end{lemma}

\begin{proof}
List the jobs $j_1,\dots,j_n$ in the order $\mathrm{DTLH}_{\max}$ processes them.
In a single window the test \eqref{eq:edf} is ``total admitted size $\le C$'' and
the comparison $s_k>s_{j_i}$ is just ``a strictly larger admitted job exists.''
Let $A_i$ be the multiset of admitted sizes after $j_1,\dots,j_i$ and
$S_i=\{s_{j_1},\dots,s_{j_i}\}$.

\emph{Invariant.} For every $i$, $A_i$ equals the set of the $k_i:=|A_i|$
smallest elements of $S_i$, where $k_i$ is the maximum cardinality of a feasible
subset of $S_i$ (equivalently, $A_i$ is the smallest-first greedy solution on
$S_i$).

The invariant holds for $i=0$. Assume it for $i-1$; write $s:=s_{j_i}$. If
$\sum A_{i-1}+s\le C$, $\mathrm{DTLH}_{\max}$ admits $j_i$. The smallest-first
greedy on $S_i$ takes the same elements: $A_{i-1}$ is the $k_{i-1}$ smallest of
$S_{i-1}$ with sum $\le C$, and appending $s$ keeps a smallest-first greedy
(greedy is order-independent), so $A_i=A_{i-1}\cup\{s\}$ is the $(k_{i-1}{+}1)$
smallest with sum $\le C$, i.e.\ the greedy optimum on $S_i$. If
$\sum A_{i-1}+s>C$ but some admitted element exceeds $s$ and removing the
\emph{largest} such element $M$ restores feasibility, then
$A_i=(A_{i-1}\setminus\{M\})\cup\{s\}$ replaces the current largest kept element
by the strictly smaller $s$; this is exactly the smallest-first greedy update
when $s$ is among the new $k_{i-1}$ smallest of $S_i$ (inserting $s$ and dropping
the largest kept element), so $A_i$ remains the greedy optimum. If no admitted
element exceeds $s$, then $s$ is at least every kept element, hence not among the
$k_{i-1}$ smallest of $S_i$, the greedy optimum is unchanged, and
$\mathrm{DTLH}_{\max}$ correctly drops $j_i$. In all cases $A_i$ is the
smallest-first greedy solution on $S_i$, which is a maximum-cardinality feasible
subset (a standard exchange argument: among feasible subsets, the smallest-first
greedy has maximum cardinality). The invariant is preserved.

By induction $A_n$ is a maximum-cardinality feasible subset of all $n$ jobs, so
$\mathrm{DTLH}_{\max}$ misses $n-|A_n|=\mathrm{Miss}_{\mathrm{OPT}}(\omega)$.
\end{proof}

\begin{remark}[The worst-case tie-break is \emph{not} optimal, already in one
window]\label{rem:onewindowbad}
Lemma~\ref{lem:onewindow} concerns the \emph{evict-largest} rule only. For the
worst-case quantity $\mathrm{Miss}_{\mathrm{DTLH}}$ of
Remark~\ref{rem:tiebreak} the conclusion \emph{fails}, even in a single window:
choosing a non-largest eligible victim wastes a slot. With capacity $C=5$
(all $a=0$, $d=5$) and sizes $3,2,1,2$ processed in that order, DTLH admits $3,2$
(total $5$); the size-$1$ arrival overflows and the worst-case tie-break evicts
the \emph{smaller} eligible job $2$ (admitted $\{3,1\}$, miss \#1); the final
size-$2$ overflows and evicts $3$ (admitted $\{1,2\}$, miss \#2). Thus the
worst-case run drops $2$ while OPT keeps $\{1,2,2\}$ and drops only the size-$3$
job, $\mathrm{Miss}_{\mathrm{OPT}}=1$. (Verified by exhaustive enumeration of all
legal tie-breaks: drop count ranges over $\{1,2\}$.) This pinpoints the source of
DTLH's loss --- the victim choice in the swap rule --- and shows it is visible
already in one window. The next section amplifies this single-window phenomenon
into an \emph{unbounded} loss.
\end{remark}

\section{Main theorem: $R=\infty$}\label{sec:main}

We resolve the headline question: DTLH is not constant-competitive.

\begin{definition}[The witness family]\label{def:family}
For an integer $r\ge 2$ let
\[
X_r=\tfrac{3r(r-1)}{2},\qquad C_r=X_r+2r,
\]
and define the single-window instance $F_r$ consisting of $r+1$ jobs, all with
arrival $a=0$ and deadline $d=C_r$ (so the only feasibility constraint is total
admitted size $\le C_r$):
\[
F_r=\Big\{\ X_r\ \Big\}\ \cup\ \big\{\,m_i:=2r+1-i\ :\ i=1,\dots,r\,\big\}.
\]
Thus $F_r$ has one ``big'' job of size $X_r$ and $r$ ``medium'' jobs of strictly
decreasing sizes $m_1=2r,\ m_2=2r-1,\ \dots,\ m_r=r+1$.
\end{definition}

We record the elementary size identities used below; all are exact integer
arithmetic.

\begin{lemma}[Arithmetic of $F_r$]\label{lem:arith}
For every $r\ge 2$:
\begin{enumerate}
\item[(i)] $\sum_{i=1}^{r} m_i = C_r$ (the medium sizes sum to exactly the
capacity);
\item[(ii)] $X_r+m_1=C_r$;
\item[(iii)] $m_i+m_{i+1}>2r$ for $i=1,\dots,r-1$, equivalently
$X_r+m_i+m_{i+1}>C_r$;
\item[(iv)] $X_r+m_i\le C_r$ for $i=1,\dots,r$;
\item[(v)] $m_i>m_{i+1}$ for $i=1,\dots,r-1$.
\end{enumerate}
\end{lemma}

\begin{proof}
$\sum_{i=1}^r m_i=\sum_{i=1}^r(2r+1-i)=r(2r+1)-\tfrac{r(r+1)}2
=\tfrac{r(3r+1)}2=\tfrac{3r(r-1)}2+2r=X_r+2r=C_r$, proving (i). For (ii),
$X_r+m_1=X_r+2r=C_r$. For (iii), $m_i+m_{i+1}=(2r+1-i)+(2r-i)=4r+1-2i$, and for
$i\le r-1$ this is $\ge 4r+1-2(r-1)=2r+3>2r$; since $C_r=X_r+2r$, adding $X_r$
gives $X_r+m_i+m_{i+1}=X_r+(m_i+m_{i+1})>X_r+2r=C_r$. For (iv),
$m_i=2r+1-i\le 2r=m_1$, so $X_r+m_i\le X_r+m_1=C_r$ by (ii). Finally (v) is
immediate from $m_i=2r+1-i$ strictly decreasing in $i$.
\end{proof}

\begin{proposition}[Offline optimum]\label{prop:family-opt}
$\mathrm{Miss}_{\mathrm{OPT}}(F_r)=1$ for every $r\ge 2$. Indeed dropping the
single big job $X_r$ leaves all $r$ medium jobs, which are feasible, and no
feasible subset of $F_r$ has more than $r$ jobs.
\end{proposition}

\begin{proof}
The medium jobs have total size $C_r$ by Lemma~\ref{lem:arith}(i), so the set
$\{m_1,\dots,m_r\}$ is feasible (total $=C_r\le C_r$); hence
$\mathrm{Miss}_{\mathrm{OPT}}(F_r)\le 1$. Conversely, $F_r$ has $r+1$ jobs of
total size $X_r+\sum m_i=X_r+C_r>C_r$, so the whole set is infeasible; any
feasible subset omits at least one job and so has at most $r$ jobs. Therefore the
maximum feasible cardinality is exactly $r$ and
$\mathrm{Miss}_{\mathrm{OPT}}(F_r)=(r+1)-r=1$.
\end{proof}

\begin{proposition}[DTLH worst-case]\label{prop:family-dtlh}
For every $r\ge 2$ there is a legal DTLH execution on $F_r$ that drops $r-1$
jobs; hence $\mathrm{Miss}_{\mathrm{DTLH}}(F_r)\ge r-1$.
\end{proposition}

\begin{proof}
Process the arrivals (all at $t=0$, a legal tie-break ordering of simultaneous
arrivals) in the order
\[
X_r,\ m_1,\ m_2,\ \dots,\ m_r,
\]
and at each swap choose the victim specified below; we verify each step satisfies
the preconditions of Algorithm~\ref{alg:dtlh}. Because all jobs share the window,
feasibility is ``total admitted size $\le C_r$'' and original size $=$ remaining
size throughout (no service precedes any admission decision).

\emph{Admit $X_r$:} total $0\le C_r$, admit; admitted $\{X_r\}$, total $X_r$.

\emph{Admit $m_1$:} total $X_r+m_1=C_r\le C_r$ by
Lemma~\ref{lem:arith}(ii), admit; admitted $\{X_r,m_1\}$, total $C_r$.

\emph{Step $i$ for $i=2,\dots,r$ (arrival $m_i$).} Inductively the admitted set
before this step is $\{X_r,m_{i-1}\}$ with total $X_r+m_{i-1}$. Adding $m_i$
gives total $X_r+m_{i-1}+m_i>C_r$ by Lemma~\ref{lem:arith}(iii)
(with index $i-1$), so the set is infeasible and line~3 fires. The eligible
victim set is
\[
L=\{k\in\{X_r,m_{i-1}\}:\ s_k>m_i\ \text{ and removing }k\text{ restores
feasibility}\}.
\]
We claim $m_{i-1}\in L$: by Lemma~\ref{lem:arith}(v), $m_{i-1}>m_i$, and after
removing $m_{i-1}$ the admitted total is $X_r+m_i\le C_r$ by
Lemma~\ref{lem:arith}(iv), so feasibility is restored. Hence $L\neq\varnothing$,
and we (legally) pick $k^\star=m_{i-1}$: DTLH drops $m_{i-1}$ and admits $m_i$.
The admitted set becomes $\{X_r,m_i\}$ with total $X_r+m_i\le C_r$, and one miss
is charged. This completes the induction for step $i$.

Over steps $i=2,\dots,r$ this drops exactly the $r-1$ distinct jobs
$m_1,m_2,\dots,m_{r-1}$. No further arrivals occur, the final admitted set
$\{X_r,m_r\}$ is feasible, and so no further miss is incurred. Thus this legal
DTLH execution drops $r-1$ jobs.
\end{proof}

\begin{theorem}[DTLH is not constant-competitive]\label{thm:main}
$\displaystyle R=\sup_{\omega:\ \mathrm{Miss}_{\mathrm{OPT}}(\omega)>0}
\frac{\mathrm{Miss}_{\mathrm{DTLH}}(\omega)}{\mathrm{Miss}_{\mathrm{OPT}}(\omega)}
=\infty.$
\end{theorem}

\begin{proof}
By Propositions~\ref{prop:family-opt} and~\ref{prop:family-dtlh}, for every
$r\ge 2$,
\[
\frac{\mathrm{Miss}_{\mathrm{DTLH}}(F_r)}{\mathrm{Miss}_{\mathrm{OPT}}(F_r)}
\ \ge\ \frac{r-1}{1}\ =\ r-1 .
\]
Letting $r\to\infty$ shows the supremum is unbounded, i.e.\ $R=\infty$.
Equivalently, no constant $c$ satisfies
$\mathrm{Miss}_{\mathrm{DTLH}}(\omega)\le c\cdot\mathrm{Miss}_{\mathrm{OPT}}(\omega)$
for all $\omega$.
\end{proof}

\begin{remark}[Mechanism]\label{rem:mechanism}
The construction isolates exactly why DTLH fails. The clairvoyant optimum removes
a \emph{single} ``linchpin'' --- the one big job $X_r$ --- after which all $r$
mediums fit perfectly (their sizes sum to the capacity). DTLH instead keeps the
big job throughout and, on each successive medium arrival, its myopic swap rule
evicts the previously admitted medium (a strictly larger feasible job) to admit
the new, slightly smaller medium. Each arrival thus wastes one medium, producing
$r-1$ drops while preserving the very job OPT would have discarded. The loss is
unbounded because the number of mediums is unbounded while the offline fix stays
a single deletion. This is the single-window slot-wasting phenomenon of
Remark~\ref{rem:onewindowbad}, scaled up.
\end{remark}

\begin{remark}[An additive lower bound, and the gap to OPT]
Since $\mathrm{Miss}_{\mathrm{OPT}}(F_r)=1$ and
$\mathrm{Miss}_{\mathrm{DTLH}}(F_r)\ge r-1$, the family also shows the
\emph{additive} gap $\mathrm{Miss}_{\mathrm{DTLH}}-\mathrm{Miss}_{\mathrm{OPT}}$
is unbounded (it is $\ge r-2$ on $F_r$), so DTLH admits neither a multiplicative
nor an additive constant guarantee against the clairvoyant optimum. By contrast
Lemma~\ref{lem:nowaste} shows it never errs on schedulable inputs.
\end{remark}

\section{The deterministic evict-largest rule}\label{sec:det}

Theorem~\ref{thm:main} concerns the underspecified policy \emph{class}
(Remark~\ref{rem:tiebreak}): the witness $F_r$ realizes $r-1$ drops only under
the adversarial victim choice (evict the smaller eligible medium). It is natural
to ask whether the unboundedness is merely a tie-break artifact. We record what
is and is not known for the canonical deterministic instantiation
\[
\mathrm{DTLH}_{\max}:\quad
\text{process simultaneous arrivals in increasing original size,}\
\text{and on every swap evict an eligible victim of maximum original size.}
\]

\begin{proposition}[$\mathrm{DTLH}_{\max}$ is single-window optimal but not
constant-competitive overall]\label{prop:det}
$\mathrm{DTLH}_{\max}$ is exactly optimal on every single common
release/deadline window (Lemma~\ref{lem:onewindow}). In particular it drops only
$1$ job on each $F_r$, so $F_r$ does \emph{not} witness unboundedness for
$\mathrm{DTLH}_{\max}$. Nonetheless $\mathrm{DTLH}_{\max}$ is not optimal across
multiple windows: there are explicit instances on which it loses a factor (a
verified instance with $\mathrm{Miss}_{\mathrm{OPT}}=1$ and
$\mathrm{Miss}_{\mathrm{DTLH}_{\max}}=5$ is given in Appendix~\ref{app:omega}).
\end{proposition}

\begin{proof}
Single-window optimality is Lemma~\ref{lem:onewindow}; applied to $F_r$ it gives
$\mathrm{Miss}_{\mathrm{DTLH}_{\max}}(F_r)=\mathrm{Miss}_{\mathrm{OPT}}(F_r)=1$.
The multi-window instance and the value $5$ are machine-verified, see
Appendix~\ref{app:omega}.
\end{proof}

\begin{remark}[Status of the deterministic ratio]\label{rem:detopen}
Theorem~\ref{thm:main} settles the headline question $R=\infty$ for DTLH as
literally stated (a rule specified only up to tie-breaks). For the single
\emph{fixed} deterministic rule $\mathrm{DTLH}_{\max}$ we have proved a lower
bound of $5$ (Appendix~\ref{app:omega}) but \emph{not} a closed-form unbounded
multi-window family; a local search over instances with
$\mathrm{Miss}_{\mathrm{OPT}}=1$ produced $\mathrm{DTLH}_{\max}$ miss counts up
to $5$ without exposing a constant ceiling, but we do not assert
$\mathrm{DTLH}_{\max}$ is non-constant-competitive as a theorem. This is the one
remaining open sub-question; it does not affect $R=\infty$ for the policy as
posed.
\end{remark}

\section{Summary}

\begin{itemize}
\item (Theorem~\ref{thm:main}) $R=\infty$: DTLH is not constant-competitive
against the clairvoyant offline optimum. The witness is the elementary
single-window family $F_r$ with $\mathrm{Miss}_{\mathrm{OPT}}=1$ and a legal DTLH
execution dropping $r-1$ jobs; the additive gap is likewise unbounded.
\item (Lemma~\ref{lem:nowaste}) DTLH is optimal whenever the input is
schedulable: $\mathrm{Miss}_{\mathrm{OPT}}=0\Rightarrow
\mathrm{Miss}_{\mathrm{DTLH}}=0$, for every tie-break.
\item (Lemma~\ref{lem:onewindow}) On any single common release/deadline window
the deterministic \emph{evict-largest} rule $\mathrm{DTLH}_{\max}$ is exactly
optimal --- but this fails for the worst-case tie-break
(Remark~\ref{rem:onewindowbad}), and it is the worst-case victim choice that
drives the unbounded loss in $F_r$.
\item (Section~\ref{sec:det}) The unbounded loss in Theorem~\ref{thm:main} uses
the adversarial tie-break; for the fixed deterministic rule
$\mathrm{DTLH}_{\max}$ we prove only a lower bound of $5$, leaving its exact
competitive ratio open.
\end{itemize}

\appendix
\section{A multi-window deterministic instance with ratio $5$}\label{app:omega}

This appendix supplies the witness used in Proposition~\ref{prop:det}; it is not
needed for the main theorem.

\begin{proposition}\label{prop:omega}
Let (triples $(a,s,d)$)
\[
\omega^\star=\big\{(8,4,21),(4,5,11),(9,1,16),(8,3,14),(5,2,16),
(1,4,16),(9,1,13),(10,1,13),(7,3,18),(8,5,25)\big\}.
\]
Then $\mathrm{Miss}_{\mathrm{OPT}}(\omega^\star)=1$, and the deterministic rule
$\mathrm{DTLH}_{\max}$ (and also a legal worst-case DTLH execution) drops $5$
jobs, so $\mathrm{Miss}_{\mathrm{DTLH}_{\max}}(\omega^\star)=5$.
\end{proposition}

\begin{proof}[Verification]
\emph{Offline optimum.} Deleting $(4,5,11)$ leaves nine jobs whose preemptive EDF
schedule meets all deadlines, e.g.
\[
\begin{array}{ll}
[1,5]\ (1,4,16) & [11,13]\ (8,3,14)\\
[5,7]\ (5,2,16) & [13,14]\ (9,1,16)\\
[7,8]\ (7,3,18) & [14,16]\ (7,3,18)\\
[8,9]\ (8,3,14) & [16,20]\ (8,4,21)\\
[9,10]\ (9,1,13)& [20,25]\ (8,5,25)\\
[10,11]\ (10,1,13)&
\end{array}
\]
(each job completes by its deadline), so $\mathrm{Miss}_{\mathrm{OPT}}\le 1$; the
full ten-job set is EDF-infeasible (a single deadline is missed, optimal by
Horn's theorem), so $\mathrm{Miss}_{\mathrm{OPT}}(\omega^\star)=1$. \emph{DTLH.}
Running $\mathrm{DTLH}_{\max}$ (advance the admitted set by EDF between epochs;
at each arrival apply Algorithm~\ref{alg:dtlh} with increasing-size arrival order
and largest-eligible victim, original-size eviction) drops the five jobs
$(4,5,11),(8,5,25),(1,4,16),(8,4,21),(8,3,14)$; this is a finite deterministic
simulation. Both computations are elementary and machine-verified.
\end{proof}

\section{Reproducibility of the machine-verified claims}\label{app:verify}

The main theorem (Theorem~\ref{thm:main}) is proved by hand in
Section~\ref{sec:main}; Lemmas~\ref{lem:arith}, \ref{lem:nowaste},
\ref{lem:onewindow} and Propositions~\ref{prop:family-opt},
\ref{prop:family-dtlh} require no computation beyond the displayed arithmetic.
For completeness we list the auxiliary numerical checks (used only for the
illustrative Remark~\ref{rem:onewindowbad}, Proposition~\ref{prop:det}, and
Appendix~\ref{app:omega}).

\paragraph{What is computed.}
\begin{itemize}
\item $\mathrm{Miss}_{\mathrm{OPT}}(\omega)$: enumerate subsets of $\omega$ in
increasing number of deletions, returning the first cardinality for which some
surviving subset is schedulable; schedulability is decided by simulating
preemptive EDF and checking all deadlines (correct by Horn's theorem).
\item $\mathrm{Miss}_{\mathrm{DTLH}}(\omega)$ (worst case): depth-first search
over all legal tie-breaks (every order of simultaneous arrivals; at every swap,
every admissible victim $k$ with $s_k>s_j$ whose removal restores feasibility),
advancing the admitted set by EDF between epochs, returning the maximum (and
minimum) number of drops.
\item $\mathrm{DTLH}_{\max}$: a single deterministic simulation (increasing-size
arrival order, largest-eligible victim, original-size eviction).
\end{itemize}

\paragraph{Outcomes (all confirmed).}
\[
\begin{array}{lccc}
\text{instance} & \mathrm{Miss}_{\mathrm{OPT}} & \min_{\text{t.b.}}\mathrm{Miss}_{\mathrm{DTLH}} & \max_{\text{t.b.}}\mathrm{Miss}_{\mathrm{DTLH}}\\\hline
F_2,\dots,F_6\ (\text{Def.~\ref{def:family}}) & 1 & 1 & 1,2,3,4,5\\
\text{one window }3,2,1,2\ (\text{Rem.~\ref{rem:onewindowbad}}) & 1 & 1 & 2\\
\omega^\star\ (\text{Prop.~\ref{prop:omega}}) & 1 & 2 & 5
\end{array}
\]
The $\max_{\text{t.b.}}$ column for $F_r$ confirms the hand proof
(Proposition~\ref{prop:family-dtlh}): the worst-case DTLH miss count on $F_r$ is
exactly $r-1$, matching the constructed execution, while
$\mathrm{Miss}_{\mathrm{OPT}}(F_r)=1$.

\end{document}
